% !TeX program = pdflatex
% Zbior zadan do czterech tomow (na podstawie niezmienionych etykiet rozdzialow).
\documentclass[11pt,a4paper,oneside]{report}
\usepackage[T1]{fontenc}
\usepackage[utf8]{inputenc}
\usepackage[polish]{babel}
\usepackage{lmodern}
\usepackage{amsmath,amssymb,amsthm,mathtools}
\usepackage{microtype}
\usepackage{geometry}
\usepackage{enumitem}
\usepackage{xcolor}
\usepackage[unicode,colorlinks=true,linkcolor=blue!45!black,urlcolor=blue!45!black]{hyperref}
\usepackage{xr-hyper}
\makeatletter
\renewcommand*\l@section{\@dottedtocline{1}{1.5em}{3.8em}}
\makeatother
\externaldocument[I-]{rachunek_prawdopodobienstwa_tom_I}[rachunek_prawdopodobienstwa_tom_I.pdf]
\externaldocument[II-]{rachunek_prawdopodobienstwa_tom_II}[rachunek_prawdopodobienstwa_tom_II.pdf]
\externaldocument[III-]{rachunek_prawdopodobienstwa_tom_III}[rachunek_prawdopodobienstwa_tom_III.pdf]
\externaldocument[IV-]{rachunek_prawdopodobienstwa_tom_IV}[rachunek_prawdopodobienstwa_tom_IV.pdf]
\geometry{margin=2.4cm}
\setlength{\parindent}{0pt}
\setlength{\parskip}{3pt}
\setlength{\emergencystretch}{2em}
\newcounter{zadanie}[section]
\renewcommand{\thesection}{\Roman{part}.\arabic{section}}
\renewcommand{\thezadanie}{\thesection.\arabic{zadanie}}
\newenvironment{zadanie}[1][]{%
  \refstepcounter{zadanie}\par\medskip\noindent\textbf{Zadanie~\thezadanie.}\if\relax\detokenize{#1}\relax\else\ \textit{(#1)}\fi\quad\ignorespaces
}{\par\medskip}
\newenvironment{wskazowka}{\par\smallskip\noindent\textit{Wskazówka.}\ \ignorespaces}{\par\smallskip}
\begin{document}
\title{Zbiór zadań do czterech tomów\\[2mm]\large Rachunek prawdopodobieństwa, procesy stochastyczne, analiza stochastyczna i statystyka matematyczna}
\author{Materiał uzupełniający do tomów I--IV}
\date{}
\maketitle
\tableofcontents
\chapter*{Jak korzystać z tego zbioru}
\addcontentsline{toc}{chapter}{Jak korzystać z tego zbioru}
Zbiór jest osobnym dokumentem i \emph{nie zmienia numeracji rozdziałów ani etykiet} czterech tomów teorii. Każdy rozdział zbioru odpowiada dokładnie jednemu rozdziałowi odpowiedniego tomu. Podano odsyłacz do teorii w formacie PDF (w czytnikach obsługujących odsyłacze między plikami).\\
Zadania oznaczone \textit{podstawowe} sprawdzają rachunek i definicje, \textit{dowodowe} wymagają uzasadnienia, \textit{praktyczne} dotyczą doboru modelu lub metody. Wskazówki podano wtedy, gdy zadanie wymaga istotnego pomysłu; nie zastępują one dowodu.\\
Aby odsyłacze do teorii działały, skompiluj najpierw tomy \texttt{rachunek\_prawdopodobienstwa\_tom\_I.tex}--\texttt{rachunek\_prawdopodobienstwa\_tom\_IV.tex}, a następnie niniejszy plik; zachowaj wszystkie pliki PDF i AUX w tym samym folderze. W czytnikach nieobsługujących odsyłaczy między PDF-ami użyj drukowanych numerów rozdziałów.

\part{Tom I. Rachunek prawdopodobieństwa}
\setcounter{section}{0}
\section{Kombinatoryka i modele skończone}
\label{zadania:I:1}
\noindent\textit{Teoria:} tom I, rozdział~\ref{I-ch:combinatorics} (\hyperref[I-ch:combinatorics]{przejdź do materiału}).\par\medskip
\begin{zadanie}[podstawowe]
Policz liczbę surjekcji ze zbioru $\{1,\ldots,7\}$ na zbiór trzyelementowy; porównaj dwa dowody.
\end{zadanie}
\begin{wskazowka}
Zastosuj zasadę włączeń i wyłączeń; alternatywnie użyj liczb Stirlinga.
\end{wskazowka}
\begin{zadanie}[dowodowe]
W urnie jest $a$ kul czerwonych, $b$ niebieskich i $c$ zielonych, gdzie $a,b,c\ge2$. Wyraź prawdopodobieństwo zobaczenia wszystkich kolorów w czterech losowaniach ze zwracaniem.
\end{zadanie}
\begin{wskazowka}
Policz dopełnienie i przecięcia zdarzeń braku koloru.
\end{wskazowka}
\begin{zadanie}[praktyczne]
Na ile sposobów można rozmieścić osiem różnych przedmiotów w trzech oznaczonych pudełkach, aby żadne nie było puste? Uogólnij na $n,k$.
\end{zadanie}
\begin{wskazowka}
Wyjaśnij różnicę między pudełkami oznaczonymi i nieoznaczonymi.
\end{wskazowka}
\section{Przestrzeń probabilistyczna i miara prawdopodobieństwa}
\label{zadania:I:2}
\noindent\textit{Teoria:} tom I, rozdział~\ref{I-ch:probability} (\hyperref[I-ch:probability]{przejdź do materiału}).\par\medskip
\begin{zadanie}[podstawowe]
Zbuduj przestrzeń probabilistyczną dla dwóch rzutów monetą o różnych prawdopodobieństwach orła; wskaż $\sigma$-ciało i miarę.
\end{zadanie}
\begin{zadanie}[dowodowe]
Podaj ciąg zdarzeń $A_n$ taki, że $\mathbb P(A_n)\to0$, ale $\mathbb P(\limsup_n A_n)=1$.
\end{zadanie}
\begin{wskazowka}
Spróbuj wykorzystać niezależność i drugi lemat Borela--Cantellego.
\end{wskazowka}
\begin{zadanie}[praktyczne]
Udowodnij ciągłość miary z góry i pokaż, dlaczego w ogólności dla miary nieskończonej potrzeba założenia skończoności pierwszego zbioru.
\end{zadanie}
\begin{wskazowka}
Rozpatrz malejące półproste w $\mathbb R$.
\end{wskazowka}
\section{Prawdopodobieństwo klasyczne, warunkowanie i niezależność}
\label{zadania:I:3}
\noindent\textit{Teoria:} tom I, rozdział~\ref{I-ch:conditioning} (\hyperref[I-ch:conditioning]{przejdź do materiału}).\par\medskip
\begin{zadanie}[podstawowe]
Dla trzech fabryk o różnych częstościach wad oblicz prawdopodobieństwo pochodzenia wadliwego produktu z każdej fabryki.
\end{zadanie}
\begin{wskazowka}
Zacznij od bezwarunkowego prawdopodobieństwa wady.
\end{wskazowka}
\begin{zadanie}[dowodowe]
Skonstruuj zdarzenia $A,B,C$ parami niezależne, lecz nie niezależne łącznie.
\end{zadanie}
\begin{wskazowka}
Wykorzystaj trzy zmienne o wartościach $\pm1$ z warunkiem na ich iloczyn.
\end{wskazowka}
\begin{zadanie}[praktyczne]
Porównaj szanse choroby przy dodatnim teście dla częstości bazowej $1\%$ i $20\%$ przy tej samej czułości i swoistości.
\end{zadanie}
\begin{wskazowka}
Wyprowadź wynik z definicji prawdopodobieństwa warunkowego.
\end{wskazowka}
\section{Zmienne losowe i rozkłady prawdopodobieństwa}
\label{zadania:I:4}
\noindent\textit{Teoria:} tom I, rozdział~\ref{I-ch:random-variables} (\hyperref[I-ch:random-variables]{przejdź do materiału}).\par\medskip
\begin{zadanie}[podstawowe]
Dla $X$ o rozkładzie jednostajnym na $(0,1)$ wyznacz rozkład $Y=-\log X$.
\end{zadanie}
\begin{wskazowka}
Najpierw oblicz dystrybuantę.
\end{wskazowka}
\begin{zadanie}[dowodowe]
Wyznacz gęstość $X/Y$ dla niezależnych dodatnich zmiennych wykładniczych o tym samym parametrze.
\end{zadanie}
\begin{wskazowka}
Użyj przekształcenia $(x,y)\mapsto(x/y,y)$ i jakobianu.
\end{wskazowka}
\begin{zadanie}[praktyczne]
Zbuduj zmienną losową o dystrybuancie ciągłej, ale bez gęstości; wyjaśnij, dlaczego ciągłość nie oznacza absolutnej ciągłości.
\end{zadanie}
\begin{wskazowka}
Wykorzystaj rozkład Cantora.
\end{wskazowka}
\section{Wartość oczekiwana, momenty i podstawowe nierówności}
\label{zadania:I:5}
\noindent\textit{Teoria:} tom I, rozdział~\ref{I-ch:expectation} (\hyperref[I-ch:expectation]{przejdź do materiału}).\par\medskip
\begin{zadanie}[podstawowe]
Udowodnij nierówność Czebyszewa przez nierówność Markowa i zbadaj przypadek równości.
\end{zadanie}
\begin{zadanie}[dowodowe]
Oblicz $\mathbb E(X^4)$ dla standardowego rozkładu normalnego dwiema metodami.
\end{zadanie}
\begin{wskazowka}
Skorzystaj z funkcji tworzącej momenty albo całkowania przez części.
\end{wskazowka}
\begin{zadanie}[praktyczne]
Dla niezależnych $X,Y$ porównaj $\operatorname{Var}(X+Y)$; skonstruuj przykład zależny, w którym wariancja sumy wynosi zero.
\end{zadanie}
\begin{wskazowka}
Wybierz $Y=-X$.
\end{wskazowka}
\section{Podstawowe rodziny rozkładów prawdopodobieństwa}
\label{zadania:I:6}
\noindent\textit{Teoria:} tom I, rozdział~\ref{I-ch:distributions} (\hyperref[I-ch:distributions]{przejdź do materiału}).\par\medskip
\begin{zadanie}[podstawowe]
Zmienna $X$ ma rozkład geometryczny na $\{1,2,\ldots\}$. Wyprowadź $\mathbb E X$ i $\operatorname{Var}X$ z szeregu geometrycznego.
\end{zadanie}
\begin{zadanie}[dowodowe]
Udowodnij granicę dwumian--Poisson przy $np_n\to\lambda>0$ dla ustalonej liczby sukcesów.
\end{zadanie}
\begin{wskazowka}
Oddziel granicę współczynnika dwumianowego i $(1-p_n)^n$.
\end{wskazowka}
\begin{zadanie}[praktyczne]
Porównaj model hipergeometryczny z dwumianowym przy losowaniu bez zwracania: wyprowadź wariancję obu liczników sukcesów.
\end{zadanie}
\begin{wskazowka}
Zapisz wskaźniki wylosowania sukcesu i ich kowariancje.
\end{wskazowka}
\section{Zbieżność zmiennych losowych i rozkładów}
\label{zadania:I:7}
\noindent\textit{Teoria:} tom I, rozdział~\ref{I-ch:convergence} (\hyperref[I-ch:convergence]{przejdź do materiału}).\par\medskip
\begin{zadanie}[podstawowe]
Pokaż, że zbieżność prawie na pewno implikuje zbieżność według prawdopodobieństwa.
\end{zadanie}
\begin{zadanie}[dowodowe]
Podaj ciąg zbieżny według prawdopodobieństwa, ale niezbieżny prawie na pewno.
\end{zadanie}
\begin{wskazowka}
Użyj tzw. maszerujących przedziałów.
\end{wskazowka}
\begin{zadanie}[praktyczne]
Podaj zmienne $X_n$ zbieżne według rozkładu do zera, lecz $\mathbb E X_n\not\to0$.
\end{zadanie}
\begin{wskazowka}
Niech rzadkie duże wartości utrzymują średnią.
\end{wskazowka}
\section{Prawa wielkich liczb}
\label{zadania:I:8}
\noindent\textit{Teoria:} tom I, rozdział~\ref{I-ch:lln} (\hyperref[I-ch:lln]{przejdź do materiału}).\par\medskip
\begin{zadanie}[podstawowe]
Udowodnij słabe prawo wielkich liczb dla niezależnych zmiennych o jednakowej średniej i ograniczonej wariancji.
\end{zadanie}
\begin{zadanie}[dowodowe]
Zbadaj, czy średnie z próby o rozkładzie Cauchy'ego zbiegają do stałej.
\end{zadanie}
\begin{wskazowka}
Oblicz funkcję charakterystyczną sumy albo skorzystaj ze stabilności.
\end{wskazowka}
\begin{zadanie}[praktyczne]
Udowodnij mocne prawo dla wskaźników Bernoulliego i wyjaśnij interpretację częstościową.
\end{zadanie}
\begin{wskazowka}
Można użyć odpowiedniego mocnego prawa z tekstu; wskaż dokładnie jego hipotezy.
\end{wskazowka}
\section{Centralne twierdzenie graniczne}
\label{zadania:I:9}
\noindent\textit{Teoria:} tom I, rozdział~\ref{I-ch:clt} (\hyperref[I-ch:clt]{przejdź do materiału}).\par\medskip
\begin{zadanie}[podstawowe]
Wyprowadź CLT dla średniej Bernoulliego i użyj jej do przybliżenia prawdopodobieństwa odsetka sukcesów.
\end{zadanie}
\begin{zadanie}[dowodowe]
Dla rozkładu dwumianowego porównaj aproksymację normalną z poprawką ciągłości i bez niej.
\end{zadanie}
\begin{wskazowka}
Rozpisz obie granice przedziału w standaryzacji.
\end{wskazowka}
\begin{zadanie}[praktyczne]
Skonstruuj przykład, w którym warunek Lindeberga nie zachodzi, i sprawdź czy granica normalna mimo to może nie wystąpić.
\end{zadanie}
\begin{wskazowka}
Rozważ pojedynczy składnik dominujący sumę.
\end{wskazowka}
\section{Warunkowa wartość oczekiwana i projekcje w $L^2$}
\label{zadania:I:10}
\noindent\textit{Teoria:} tom I, rozdział~\ref{I-ch:condexp} (\hyperref[I-ch:condexp]{przejdź do materiału}).\par\medskip
\begin{zadanie}[podstawowe]
Dla $X,Y$ o skończonych drugich momentach udowodnij, że $\mathbb E(X\mid\sigma(Y))$ jest rzutem ortogonalnym w $L^2$.
\end{zadanie}
\begin{zadanie}[dowodowe]
Oblicz $\mathbb E(X\mid X+Y)$ dla niezależnych standardowych zmiennych normalnych.
\end{zadanie}
\begin{wskazowka}
Użyj symetrii i własności warunkowej wartości oczekiwanej.
\end{wskazowka}
\begin{zadanie}[praktyczne]
Pokaż na przykładzie, że $\mathbb E(X\mid\mathcal G)=\mathbb E X$ nie implikuje niezależności $X$ i $\mathcal G$.
\end{zadanie}
\begin{wskazowka}
Poszukaj zależności nieliniowej o zerowej średniej warunkowej.
\end{wskazowka}
\section{Dalsze twierdzenia graniczne i koncentracja prawdopodobieństwa}
\label{zadania:I:11}
\noindent\textit{Teoria:} tom I, rozdział~\ref{I-ch:advancedlimits} (\hyperref[I-ch:advancedlimits]{przejdź do materiału}).\par\medskip
\begin{zadanie}[podstawowe]
Udowodnij nierówność Hoeffdinga dla średniej niezależnych zmiennych Bernoulliego w oparciu o lemat wykładniczy.
\end{zadanie}
\begin{wskazowka}
Zastosuj wykładniczą nierówność Markowa.
\end{wskazowka}
\begin{zadanie}[dowodowe]
Porównaj oszacowania Czebyszewa i Hoeffdinga dla $n=100$ i odchylenia średniej o $0{,}1$.
\end{zadanie}
\begin{zadanie}[praktyczne]
Wykaż, że zbieżność według prawdopodobieństwa oraz jednostajna całkowalność implikują zbieżność w $L^1$ przy odpowiedniej wartości granicznej.
\end{zadanie}
\begin{wskazowka}
Rozdziel całkę na część ograniczoną i ogon.
\end{wskazowka}
\section{Łańcuchy Markowa w czasie dyskretnym}
\label{zadania:I:12}
\noindent\textit{Teoria:} tom I, rozdział~\ref{I-ch:markov} (\hyperref[I-ch:markov]{przejdź do materiału}).\par\medskip
\begin{zadanie}[podstawowe]
Dla łańcucha o macierzy $P=\left(\begin{smallmatrix}0{,}8&0{,}2\\0{,}3&0{,}7\end{smallmatrix}\right)$ znajdź rozkład stacjonarny i tempo zbieżności.
\end{zadanie}
\begin{zadanie}[dowodowe]
Zbuduj łańcuch nieredukowalny, lecz okresowy, który nie zbiega do rozkładu stacjonarnego od każdej wartości początkowej.
\end{zadanie}
\begin{wskazowka}
Rozważ ruch naprzemienny na dwóch stanach.
\end{wskazowka}
\begin{zadanie}[praktyczne]
Wyznacz średni czas pierwszego osiągnięcia stanu pochłaniającego w łańcuchu trójstanowym.
\end{zadanie}
\begin{wskazowka}
Ułóż układ równań z krokiem wstecz.
\end{wskazowka}
\section{Filtracje, adaptacja i czasy zatrzymania}
\label{zadania:I:13}
\noindent\textit{Teoria:} tom I, rozdział~\ref{I-ch:filtracje} (\hyperref[I-ch:filtracje]{przejdź do materiału}).\par\medskip
\begin{zadanie}[podstawowe]
Wykaż, że pierwszy czas osiągnięcia zbioru przez adaptowany proces dyskretny jest czasem zatrzymania.
\end{zadanie}
\begin{zadanie}[dowodowe]
Podaj przykład zmiennej losowej zależnej od przyszłych obserwacji, która nie jest czasem zatrzymania.
\end{zadanie}
\begin{zadanie}[praktyczne]
Dla symetrycznego spaceru losowego zbadaj ograniczone czasy wyjścia z przedziału i zastosuj zatrzymywanie.
\end{zadanie}
\begin{wskazowka}
Najpierw zatrzymaj proces w $\tau\wedge n$.
\end{wskazowka}
\section{Martyngały: rachunek, strategie i nierówności}
\label{zadania:I:14}
\noindent\textit{Teoria:} tom I, rozdział~\ref{I-ch:martyngaly} (\hyperref[I-ch:martyngaly]{przejdź do materiału}).\par\medskip
\begin{zadanie}[podstawowe]
Udowodnij martyngałowość sum scentralizowanych niezależnych zmiennych całkowalnych.
\end{zadanie}
\begin{zadanie}[dowodowe]
Dla spaceru losowego rozwiąż problem ruiny gracza przy skończonych barierach za pomocą martyngału.
\end{zadanie}
\begin{wskazowka}
Dobierz martyngał liniowy albo wykładniczy zależnie od dryfu.
\end{wskazowka}
\begin{zadanie}[praktyczne]
Zbuduj strategię obstawiania, która nie łamie zasady uczciwej gry pomimo adaptacji.
\end{zadanie}
\begin{wskazowka}
Sprawdź mierzalność przewidywalną współczynników.
\end{wskazowka}
\section{Podmartyngały, nadmartyngały i twierdzenia o zbieżności}
\label{zadania:I:15}
\noindent\textit{Teoria:} tom I, rozdział~\ref{I-ch:nadmartyngaly} (\hyperref[I-ch:nadmartyngaly]{przejdź do materiału}).\par\medskip
\begin{zadanie}[podstawowe]
Pokaż, że kwadrat martyngału o skończonym drugim momencie jest podmartyngałem.
\end{zadanie}
\begin{zadanie}[dowodowe]
Wyprowadź nierówność maksymalną Dooba dla nieujemnego podmartyngału w czasie dyskretnym.
\end{zadanie}
\begin{wskazowka}
Rozważ pierwszy czas przekroczenia poziomu.
\end{wskazowka}
\begin{zadanie}[praktyczne]
Podaj przykład martyngału zbieżnego prawie na pewno, ale niezbieżnego w $L^1$; wyjaśnij rolę jednostajnej całkowalności.
\end{zadanie}
\begin{wskazowka}
Użyj krytycznego procesu mnożnikowego albo zatrzymanego spaceru.
\end{wskazowka}


\part{Tom II. Procesy stochastyczne}
\setcounter{section}{0}
\section{Procesy stochastyczne: od rozkładów do trajektorii}
\label{zadania:II:1}
\noindent\textit{Teoria:} tom II, rozdział~\ref{II-ch:procesy} (\hyperref[II-ch:procesy]{przejdź do materiału}).\par\medskip
\begin{zadanie}[podstawowe]
Zbuduj proces z niezależnych rzutów monetą i określ jego rozkłady skończeniewymiarowe.
\end{zadanie}
\begin{zadanie}[dowodowe]
Pokaż, że równość wszystkich rozkładów skończeniewymiarowych nie gwarantuje nierozróżnialności dwóch wersji procesu.
\end{zadanie}
\begin{wskazowka}
Użyj modyfikacji w losowej chwili.
\end{wskazowka}
\begin{zadanie}[praktyczne]
Dla procesu $X_t=Zt$ sprawdź stacjonarność i niezależność przyrostów.
\end{zadanie}
\section{Konstrukcja procesów i regularność trajektorii}
\label{zadania:II:2}
\noindent\textit{Teoria:} tom II, rozdział~\ref{II-ch:konstrukcja} (\hyperref[II-ch:konstrukcja]{przejdź do materiału}).\par\medskip
\begin{zadanie}[podstawowe]
Dla kowariancji $K(s,t)=\min(s,t)$ sprawdź dodatnią półokreśloność każdej macierzy kowariancji.
\end{zadanie}
\begin{wskazowka}
Zinterpretuj ją jako iloczyn skalarny wskaźników przedziałów.
\end{wskazowka}
\begin{zadanie}[dowodowe]
Wyprowadź warunek momentowy gwarantujący ciągłą modyfikację z kryterium Kołmogorowa.
\end{zadanie}
\begin{zadanie}[praktyczne]
Podaj proces ze zgodnymi rozkładami skończeniewymiarowymi, ale z trudnością dotyczącą wyboru wersji mierzalnej względem wspólnej zmiennej $(t,\omega)$.
\end{zadanie}
\begin{wskazowka}
Rozróżnij istnienie procesu i regularność trajektorii.
\end{wskazowka}
\section{Filtracje i obserwacja procesów w czasie ciągłym}
\label{zadania:II:3}
\noindent\textit{Teoria:} tom II, rozdział~\ref{II-ch:filtracje} (\hyperref[II-ch:filtracje]{przejdź do materiału}).\par\medskip
\begin{zadanie}[podstawowe]
Sprawdź prawostronną ciągłość naturalnej filtracji procesu Wienera po standardowym uzupełnieniu.
\end{zadanie}
\begin{wskazowka}
Ustal, jakie twierdzenie z książki jest tu wykorzystywane.
\end{wskazowka}
\begin{zadanie}[dowodowe]
Dla procesu skokowego określ filtrację naturalną i filtrację zawierającą informację o przyszłym skoku.
\end{zadanie}
\begin{zadanie}[praktyczne]
Wykaż, że czas pierwszego trafienia przez ciągły proces adaptowany do zbioru domkniętego jest czasem zatrzymania.
\end{zadanie}
\begin{wskazowka}
Zapisz zdarzenie $\{\tau\le t\}$ przez przybliżenia wymierne.
\end{wskazowka}
\section{Proces Poissona: losowe zliczanie zdarzeń}
\label{zadania:II:4}
\noindent\textit{Teoria:} tom II, rozdział~\ref{II-ch:poisson-process} (\hyperref[II-ch:poisson-process]{przejdź do materiału}).\par\medskip
\begin{zadanie}[podstawowe]
Udowodnij rozkład odstępów czasowych procesu Poissona.
\end{zadanie}
\begin{zadanie}[dowodowe]
Przy warunku $N_T=n$ znajdź rozkład chwil skoków i wyznacz średnią pierwszej chwili zdarzenia.
\end{zadanie}
\begin{wskazowka}
Użyj statystyk porządkowych $n$ zmiennych jednostajnych.
\end{wskazowka}
\begin{zadanie}[praktyczne]
Dla niejednorodnego procesu Poissona z intensywnością $\lambda(t)$ oblicz $\mathbb E N_T$ i rozkład $N_T$.
\end{zadanie}
\begin{wskazowka}
Zastosuj zmianę czasu.
\end{wskazowka}
\section{Łańcuchy Markowa w czasie ciągłym}
\label{zadania:II:5}
\noindent\textit{Teoria:} tom II, rozdział~\ref{II-ch:ctmc} (\hyperref[II-ch:ctmc]{przejdź do materiału}).\par\medskip
\begin{zadanie}[podstawowe]
Dla dwustanowego łańcucha o intensywnościach $a,b>0$ wyprowadź macierz przejścia $P_t$.
\end{zadanie}
\begin{wskazowka}
Rozwiąż równania Kołmogorowa.
\end{wskazowka}
\begin{zadanie}[dowodowe]
Pokaż, że proces o intensywności wyjścia $\lambda_n=n^2$ może eksplodować w skończonym czasie.
\end{zadanie}
\begin{wskazowka}
Zbadaj sumę niezależnych czasów wykładniczych.
\end{wskazowka}
\begin{zadanie}[praktyczne]
Wyznacz rozkład stacjonarny skończonego nieodwracalnego łańcucha czasu ciągłego i sprawdź równania równowagi.
\end{zadanie}
\section{Procesy gaussowskie: od kowariancji do prawa procesu}
\label{zadania:II:6}
\noindent\textit{Teoria:} tom II, rozdział~\ref{II-ch:gauss} (\hyperref[II-ch:gauss]{przejdź do materiału}).\par\medskip
\begin{zadanie}[podstawowe]
Pokaż, że proces gaussowski wyznaczony przez średnią i kowariancję ma jednoznaczne rozkłady skończeniewymiarowe.
\end{zadanie}
\begin{zadanie}[dowodowe]
Dla $X_t=Z\cos t+Y\sin t$ z niezależnymi standardowymi normalnymi wyznacz kowariancję i zbadaj stacjonarność.
\end{zadanie}
\begin{zadanie}[praktyczne]
Sprawdź dodatnią półokreśloność funkcji $K(s,t)=e^{-|t-s|}$.
\end{zadanie}
\begin{wskazowka}
Skorzystaj z reprezentacji kowariancji procesu Ornsteina--Uhlenbecka.
\end{wskazowka}
\section{Proces Wienera: własności probabilistyczne i geometria trajektorii}
\label{zadania:II:7}
\noindent\textit{Teoria:} tom II, rozdział~\ref{II-ch:wiener} (\hyperref[II-ch:wiener]{przejdź do materiału}).\par\medskip
\begin{zadanie}[podstawowe]
Wyprowadź zasadę odbicia Browna i oblicz rozkład $\sup_{0\le s\le T}W_s$.
\end{zadanie}
\begin{zadanie}[dowodowe]
Udowodnij, że trajektorie Wienera mają niemal na pewno nieskończoną wariację na każdym przedziale.
\end{zadanie}
\begin{wskazowka}
Porównaj z wariacją kwadratową.
\end{wskazowka}
\begin{zadanie}[praktyczne]
Znajdź rozkład pierwszej chwili osiągnięcia poziomu $a>0$ przez Browna.
\end{zadanie}
\begin{wskazowka}
Zróżniczkuj dystrybuantę wyprowadzoną z odbicia.
\end{wskazowka}
\section{Ruch Browna w wielu wymiarach i dalsze własności trajektorii}
\label{zadania:II:8}
\noindent\textit{Teoria:} tom II, rozdział~\ref{II-ch:brown-multi} (\hyperref[II-ch:brown-multi]{przejdź do materiału}).\par\medskip
\begin{zadanie}[podstawowe]
Dla $d$-wymiarowego Browna i $t>0$ oblicz rozkład $\|B_t\|^2/t$.
\end{zadanie}
\begin{zadanie}[dowodowe]
Zbadaj niezależność projekcji Browna na dwa ortogonalne podprzestrzenie.
\end{zadanie}
\begin{zadanie}[praktyczne]
Wyznacz prawdopodobieństwo opuszczenia kuli przez Browna w zadanym czasie przez oszacowanie maksymalnego odchylenia.
\end{zadanie}
\begin{wskazowka}
Użyj odpowiedniej nierówności maksymalnej współrzędnych.
\end{wskazowka}
\section{Procesy odnowy: zliczanie zdarzeń bez własności braku pamięci}
\label{zadania:II:9}
\noindent\textit{Teoria:} tom II, rozdział~\ref{II-ch:renewal} (\hyperref[II-ch:renewal]{przejdź do materiału}).\par\medskip
\begin{zadanie}[podstawowe]
Dla procesu odnowy z czasami deterministycznymi oblicz $N_t$ i pokaż różnicę wobec procesu Poissona.
\end{zadanie}
\begin{zadanie}[dowodowe]
Wyprowadź równanie odnowy dla średniej liczby zdarzeń.
\end{zadanie}
\begin{wskazowka}
Warunkuj względem pierwszego odstępu.
\end{wskazowka}
\begin{zadanie}[praktyczne]
Przy czasach o rozkładzie Erlanga wyjaśnij, dlaczego proces zliczający nie ma niezależnych przyrostów.
\end{zadanie}
\section{Procesy rozgałęziające: losowy wzrost i wymieranie populacji}
\label{zadania:II:10}
\noindent\textit{Teoria:} tom II, rozdział~\ref{II-ch:branching} (\hyperref[II-ch:branching]{przejdź do materiału}).\par\medskip
\begin{zadanie}[podstawowe]
W procesie Galtona--Watsona o rozkładzie potomstwa $\mathbb P(0)=1/4$, $\mathbb P(2)=3/4$ wyznacz prawdopodobieństwo wymarcia.
\end{zadanie}
\begin{wskazowka}
Rozwiąż równanie punktu stałego funkcji tworzącej.
\end{wskazowka}
\begin{zadanie}[dowodowe]
Dla procesu Yule'a wyprowadź $\mathbb E N_t$ z równań intensywności.
\end{zadanie}
\begin{zadanie}[praktyczne]
Pokaż, że przy średniej liczbie potomstwa równej jeden proces może wymrzeć prawie na pewno mimo stałej wartości oczekiwanej populacji.
\end{zadanie}
\section{Martyngały w czasie ciągłym: informacja, regularność i lokalizacja}
\label{zadania:II:11}
\noindent\textit{Teoria:} tom II, rozdział~\ref{II-ch:mart-cont} (\hyperref[II-ch:mart-cont]{przejdź do materiału}).\par\medskip
\begin{zadanie}[podstawowe]
Udowodnij, że ciągły martyngał lokalny nie musi być martyngałem właściwym.
\end{zadanie}
\begin{wskazowka}
Poszukaj przykładu odwrotności wymiarowego procesu Bessela.
\end{wskazowka}
\begin{zadanie}[dowodowe]
Dla całkowalnego martyngału oblicz $\mathbb E(M_t\mid\mathcal F_s)$ i omów zależność od filtracji.
\end{zadanie}
\begin{zadanie}[praktyczne]
Skonstruuj ciąg czasów zatrzymania lokalizujących lokalny martyngał i wskaż, dlaczego potrzebna jest całkowalność zatrzymań.
\end{zadanie}
\section{Zatrzymywanie, nierówności maksymalne i zbieżność martyngałów}
\label{zadania:II:12}
\noindent\textit{Teoria:} tom II, rozdział~\ref{II-ch:stop-cont} (\hyperref[II-ch:stop-cont]{przejdź do materiału}).\par\medskip
\begin{zadanie}[podstawowe]
Udowodnij nierówność Dooba w $L^2$ dla martyngału ciągłego ograniczonego w $L^2$.
\end{zadanie}
\begin{zadanie}[dowodowe]
Sprawdź warunki, w których $\mathbb E M_\tau=\mathbb E M_0$ dla nieograniczonego czasu zatrzymania.
\end{zadanie}
\begin{wskazowka}
Nie wystarczy sama skończoność $\tau$ prawie na pewno.
\end{wskazowka}
\begin{zadanie}[praktyczne]
Podaj przykład pokazujący, że granica lokalizacji nie pozwala zawsze wymienić granicy z wartością oczekiwaną.
\end{zadanie}
\section{Wariacja kwadratowa i kowariacja ciągłych martyngałów}
\label{zadania:II:13}
\noindent\textit{Teoria:} tom II, rozdział~\ref{II-ch:quadratic-cont} (\hyperref[II-ch:quadratic-cont]{przejdź do materiału}).\par\medskip
\begin{zadanie}[podstawowe]
Dla Browna oblicz granicę sum kwadratów przyrostów na podziałach jednostajnych.
\end{zadanie}
\begin{zadanie}[dowodowe]
Udowodnij polaryzacyjną tożsamość $[M,N]=([M+N]-[M]-[N])/2$.
\end{zadanie}
\begin{zadanie}[praktyczne]
Wyjaśnij różnicę między $[M]$ a przewidywalnym nawiasem $\langle M\rangle$ i podaj przykład, w którym są różne.
\end{zadanie}
\begin{wskazowka}
Wybierz martyngał ze skokami.
\end{wskazowka}
\section{Martyngały ze skokami, kompensacja i granica zakresu tomu II}
\label{zadania:II:14}
\noindent\textit{Teoria:} tom II, rozdział~\ref{II-ch:jump-mart} (\hyperref[II-ch:jump-mart]{przejdź do materiału}).\par\medskip
\begin{zadanie}[podstawowe]
Dla procesu Poissona pokaż, że $M_t=N_t-\lambda t$ jest martyngałem i oblicz $[M]_t$ oraz $\langle M\rangle_t$.
\end{zadanie}
\begin{zadanie}[dowodowe]
Wyprowadź kompensator zliczania z intensywnością deterministyczną $\lambda(t)$.
\end{zadanie}
\begin{zadanie}[praktyczne]
Zbadaj różnicę pomiędzy wariacją kwadratową procesu ze skokami i jego ciągłą częścią martyngałową.
\end{zadanie}


\part{Tom III. Analiza stochastyczna}
\setcounter{section}{0}
\section{Całka Wienera dla deterministycznych funkcji}
\label{zadania:III:1}
\noindent\textit{Teoria:} tom III, rozdział~\ref{III-ch:wiener-integral} (\hyperref[III-ch:wiener-integral]{przejdź do materiału}).\par\medskip
\begin{zadanie}[podstawowe]
Udowodnij izometrię całki Wienera dla funkcji schodkowych.
\end{zadanie}
\begin{zadanie}[dowodowe]
Oblicz rozkład $\int_{[0,T]} t\,\mathrm dW_t$ i jego wariancję.
\end{zadanie}
\begin{zadanie}[praktyczne]
Wykaż, że $\int_{[0,T]} f(t)\,\mathrm dW_t$ i $\int_{[0,T]} g(t)\,\mathrm dW_t$ są niezależne przy $\langle f,g\rangle_{L^2}=0$.
\end{zadanie}
\begin{wskazowka}
Wykorzystaj wspólny rozkład gaussowski.
\end{wskazowka}
\section{Całka Itô: losowe funkcje podcałkowe}
\label{zadania:III:2}
\noindent\textit{Teoria:} tom III, rozdział~\ref{III-ch:ito-construction} (\hyperref[III-ch:ito-construction]{przejdź do materiału}).\par\medskip
\begin{zadanie}[podstawowe]
Zdefiniuj całkę Itô dla prostego procesu przewidywalnego i wyprowadź izometrię.
\end{zadanie}
\begin{zadanie}[dowodowe]
Oblicz $\mathbb E(\int_{[0,T]}W_t\,\mathrm dW_t)^2$.
\end{zadanie}
\begin{wskazowka}
Nie traktuj całki jako całki Riemanna--Stieltjesa.
\end{wskazowka}
\begin{zadanie}[praktyczne]
Podaj przykład procesu adaptowanego, który nie należy do klasy procesów całkowalnych kwadratowo na danym horyzoncie.
\end{zadanie}
\section{Wzór Itô: rachunek różniczkowy z wariacją kwadratową}
\label{zadania:III:3}
\noindent\textit{Teoria:} tom III, rozdział~\ref{III-ch:ito-formula} (\hyperref[III-ch:ito-formula]{przejdź do materiału}).\par\medskip
\begin{zadanie}[podstawowe]
Zastosuj wzór Itô do $f(x)=x^4$ i wyznacz $\mathbb E W_T^4$.
\end{zadanie}
\begin{zadanie}[dowodowe]
Oblicz $\mathrm d(e^{t/2}\cos W_t)$ za pomocą wzoru Itô.
\end{zadanie}
\begin{zadanie}[praktyczne]
Wyprowadź równanie dla $\log X_t$ przy geometrycznym ruchu Browna i wskaż poprawkę dryfu.
\end{zadanie}
\begin{wskazowka}
Pamiętaj o drugim rzędzie.
\end{wskazowka}
\section{Całka Stratonovicha i jej związek z całką Itô}
\label{zadania:III:4}
\noindent\textit{Teoria:} tom III, rozdział~\ref{III-ch:stratonovich} (\hyperref[III-ch:stratonovich]{przejdź do materiału}).\par\medskip
\begin{zadanie}[podstawowe]
Oblicz obie całki $\int_{[0,T]}W_t\circ\mathrm dW_t$ i $\int_{[0,T]}W_t\,\mathrm dW_t$.
\end{zadanie}
\begin{zadanie}[dowodowe]
Wyprowadź regułę łańcuchową Stratonovicha dla $f(W_t)$.
\end{zadanie}
\begin{zadanie}[praktyczne]
Dla SDE $\mathrm dX_t=b(X_t)\,\mathrm dt+\sigma(X_t)\circ\mathrm dW_t$ wyprowadź dryf równania Itô.
\end{zadanie}
\begin{wskazowka}
Oblicz kowariację odpowiedniego integrandu z Browniem.
\end{wskazowka}
\section{Stochastyczne równania różniczkowe: konstrukcja rozwiązań}
\label{zadania:III:5}
\noindent\textit{Teoria:} tom III, rozdział~\ref{III-ch:sde-existence} (\hyperref[III-ch:sde-existence]{przejdź do materiału}).\par\medskip
\begin{zadanie}[podstawowe]
Udowodnij jednoznaczność rozwiązań SDE z globalnie Lipschitzowskimi współczynnikami.
\end{zadanie}
\begin{wskazowka}
Zastosuj oszacowanie BDG lub izometrię i nierówność Gronwalla.
\end{wskazowka}
\begin{zadanie}[dowodowe]
Rozwiąż jawnie równanie $\mathrm dX_t=aX_t\,\mathrm dt+bX_t\,\mathrm dW_t$.
\end{zadanie}
\begin{zadanie}[praktyczne]
Podaj przykład współczynnika tylko lokalnie Lipschitzowskiego, dla którego możliwość eksplozji wymaga osobnej analizy.
\end{zadanie}
\section{Własności rozwiązań SDE, procesy Markowa i aproksymacje}
\label{zadania:III:6}
\noindent\textit{Teoria:} tom III, rozdział~\ref{III-ch:sde-properties} (\hyperref[III-ch:sde-properties]{przejdź do materiału}).\par\medskip
\begin{zadanie}[podstawowe]
Wyprowadź generator dyfuzji jednowymiarowej z jej równania SDE.
\end{zadanie}
\begin{zadanie}[dowodowe]
Dla schematu Eulera--Maruyamy przeprowadź symulację geometrycznego Browna i porównaj błąd dla kroków $h$ i $h/2$.
\end{zadanie}
\begin{zadanie}[praktyczne]
Wykaż własność Markowa rozwiązania przy odpowiednich warunkach jednoznaczności i niezależności przyrostów.
\end{zadanie}
\begin{wskazowka}
Rozpocznij proces w chwili $s$ od deterministycznego stanu.
\end{wskazowka}
\section{Problem martyngałowy i twierdzenie Feynmana--Kaca}
\label{zadania:III:7}
\noindent\textit{Teoria:} tom III, rozdział~\ref{III-ch:martingale-fk} (\hyperref[III-ch:martingale-fk]{przejdź do materiału}).\par\medskip
\begin{zadanie}[podstawowe]
Wyprowadź równanie wsteczne Kołmogorowa dla oczekiwań funkcjonałów dyfuzji.
\end{zadanie}
\begin{zadanie}[dowodowe]
Zastosuj Feynmana--Kaca do równania ciepła z potencjałem stałym.
\end{zadanie}
\begin{zadanie}[praktyczne]
Sprawdź problem martyngałowy dla rozwiązania SDE o stałych współczynnikach.
\end{zadanie}
\begin{wskazowka}
Zastosuj wzór Itô do funkcji testowej.
\end{wskazowka}
\section{Zmiana miary probabilistycznej i twierdzenie Girsanowa}
\label{zadania:III:8}
\noindent\textit{Teoria:} tom III, rozdział~\ref{III-ch:girsanov} (\hyperref[III-ch:girsanov]{przejdź do materiału}).\par\medskip
\begin{zadanie}[podstawowe]
Wyprowadź gęstość Radona--Nikodyma dla Browna z deterministycznym dryfem i sprawdź, że jej średnia wynosi jeden.
\end{zadanie}
\begin{zadanie}[dowodowe]
Zastosuj Girsanowa do wyceny oczekiwania zależnego od geometrycznego Browna po zmianie dryfu.
\end{zadanie}
\begin{zadanie}[praktyczne]
Wyjaśnij, dlaczego lokalny wykładnik stochastyczny nie musi być prawdziwym martyngałem.
\end{zadanie}
\begin{wskazowka}
Podaj warunek wystarczający Novikova i rozważ granice jego zastosowania.
\end{wskazowka}
\section{Miara Wienera i twierdzenie Camerona--Martina}
\label{zadania:III:9}
\noindent\textit{Teoria:} tom III, rozdział~\ref{III-ch:cameron-martin} (\hyperref[III-ch:cameron-martin]{przejdź do materiału}).\par\medskip
\begin{zadanie}[podstawowe]
Oblicz normę Camerona--Martina funkcji $h(t)=ct$ na $[0,T]$.
\end{zadanie}
\begin{zadanie}[dowodowe]
Pokaż, że przesunięcie Browna o funkcję o skoku prowadzi do miar wzajemnie singularnych.
\end{zadanie}
\begin{wskazowka}
Zauważ, że nośnik miary Wienera leży w przestrzeni trajektorii ciągłych.
\end{wskazowka}
\begin{zadanie}[praktyczne]
Wyprowadź wzór gęstości zmiany miary dla przesunięcia absolutnie ciągłego o pochodnej w $L^2$.
\end{zadanie}
\section{Procesy Lévy'ego i ich wykładniki charakterystyczne}
\label{zadania:III:10}
\noindent\textit{Teoria:} tom III, rozdział~\ref{III-ch:levy} (\hyperref[III-ch:levy]{przejdź do materiału}).\par\medskip
\begin{zadanie}[podstawowe]
Sprawdź wykładnik charakterystyczny złożonego procesu Poissona.
\end{zadanie}
\begin{zadanie}[dowodowe]
Wykaż, że rozkład procesu Lévy'ego jest nieskończenie podzielny dla każdej chwili $t$.
\end{zadanie}
\begin{zadanie}[praktyczne]
Dla procesu stabilnego o indeksie $\alpha\in(0,2)$ porównaj istnienie wariancji z przypadkiem Browna.
\end{zadanie}
\begin{wskazowka}
Zbadaj zachowanie ciężkich ogonów.
\end{wskazowka}
\section{Losowe miary skoków i całkowanie względem miar Poissona}
\label{zadania:III:11}
\noindent\textit{Teoria:} tom III, rozdział~\ref{III-ch:poisson-random-measure} (\hyperref[III-ch:poisson-random-measure]{przejdź do materiału}).\par\medskip
\begin{zadanie}[podstawowe]
Dla całki względem skompensowanej miary Poissona wyprowadź izometrię dla prostych integrandów.
\end{zadanie}
\begin{zadanie}[dowodowe]
Oblicz wariancję złożonego procesu Poissona o skokach o skończonym drugim momencie.
\end{zadanie}
\begin{zadanie}[praktyczne]
Wyjaśnij różnicę między kompensacją małych i dużych skoków przy całkowalności skończonego pierwszego momentu.
\end{zadanie}
\section{Półmartyngały i wzór Itô dla procesów ze skokami}
\label{zadania:III:12}
\noindent\textit{Teoria:} tom III, rozdział~\ref{III-ch:semimartingale-jumps} (\hyperref[III-ch:semimartingale-jumps]{przejdź do materiału}).\par\medskip
\begin{zadanie}[podstawowe]
Wyprowadź wzór Itô dla $f(X_t)=X_t^2$ przy semimartyngale ze skokami.
\end{zadanie}
\begin{zadanie}[dowodowe]
Dla skompensowanego procesu Poissona oblicz wariację kwadratową i przewidywalny nawias.
\end{zadanie}
\begin{zadanie}[praktyczne]
Pokaż, dlaczego przy skokach we wzorze Itô konieczny jest składnik $f(X_s)-f(X_{s-})-f'(X_{s-})\Delta X_s$.
\end{zadanie}
\section{Stochastyczne równania różniczkowe ze skokami}
\label{zadania:III:13}
\noindent\textit{Teoria:} tom III, rozdział~\ref{III-ch:jump-sde} (\hyperref[III-ch:jump-sde]{przejdź do materiału}).\par\medskip
\begin{zadanie}[podstawowe]
Rozwiąż liniowe SDE napędzane procesem Poissona, opisując zachowanie pomiędzy skokami i w chwilach skoków.
\end{zadanie}
\begin{zadanie}[dowodowe]
Wykaż jednoznaczność rozwiązania SDE ze skokami przy odpowiednich warunkach Lipschitza i całkowalności.
\end{zadanie}
\begin{zadanie}[praktyczne]
Dla procesu z intensywnością $\lambda$ i skokami stałej wielkości porównaj symulację dokładną oraz dyskretyzację czasową.
\end{zadanie}
\section{Rozwiązania silne i słabe oraz twierdzenie Yamada--Watanabego}
\label{zadania:III:14}
\noindent\textit{Teoria:} tom III, rozdział~\ref{III-ch:weak-strong-sde} (\hyperref[III-ch:weak-strong-sde]{przejdź do materiału}).\par\medskip
\begin{zadanie}[podstawowe]
Wyjaśnij różnicę między rozwiązaniem silnym a słabym dla ustalonego SDE.
\end{zadanie}
\begin{zadanie}[dowodowe]
Udowodnij implikację jednoznaczności trajektorialnej i istnienia słabego rozwiązania do istnienia silnego w sytuacji objętej twierdzeniem Yamada--Watanabego.
\end{zadanie}
\begin{wskazowka}
Zidentyfikuj dokładnie wykorzystywane twierdzenie i jego założenia.
\end{wskazowka}
\begin{zadanie}[praktyczne]
Podaj przykład, w którym jednoznaczność rozkładu nie implikuje jednoznaczności trajektorialnej.
\end{zadanie}
\section{Całki po trajektoriach i funkcjonały miary Wienera}
\label{zadania:III:15}
\noindent\textit{Teoria:} tom III, rozdział~\ref{III-ch:path-integrals} (\hyperref[III-ch:path-integrals]{przejdź do materiału}).\par\medskip
\begin{zadanie}[podstawowe]
Wyprowadź kowariancję mostu Browna z jego reprezentacji liniowej.
\end{zadanie}
\begin{zadanie}[dowodowe]
Oblicz wariancję $\int_{[0,T]}B_t\,\mathrm dt$ dla mostu od zera do zera.
\end{zadanie}
\begin{zadanie}[praktyczne]
Wyznacz wyrażenie Feynmana--Kaca dla potencjału kwadratowego i zinterpretuj transformację Laplace'a funkcjonału trajektorii.
\end{zadanie}
\section{Całki Feynmana, obrót Wicka i mechanika kwantowa}
\label{zadania:III:16}
\noindent\textit{Teoria:} tom III, rozdział~\ref{III-ch:feynman-path} (\hyperref[III-ch:feynman-path]{przejdź do materiału}).\par\medskip
\begin{zadanie}[podstawowe]
Sprawdź wzór wolnego propagatora Schrödingera przez transformatę Fouriera.
\end{zadanie}
\begin{zadanie}[dowodowe]
Wyprowadź ślad półgrupy oscylatora harmonicznego z jego widma i porównaj z diagonalną całką jądra Mehlera.
\end{zadanie}
\begin{zadanie}[praktyczne]
Wyjaśnij, dlaczego oscylacyjnej całki Feynmana nie można utożsamiać z całką względem skończonej miary zespolonej na przestrzeni trajektorii.
\end{zadanie}
\begin{wskazowka}
Porównaj konstrukcję euklidesową i unitarną.
\end{wskazowka}


\part{Tom IV. Statystyka matematyczna}
\setcounter{section}{0}
\section{Eksperyment statystyczny i model danych}
\label{zadania:IV:1}
\noindent\textit{Teoria:} tom IV, rozdział~\ref{IV-ch:eksperyment} (\hyperref[IV-ch:eksperyment]{przejdź do materiału}).\par\medskip
\begin{zadanie}[podstawowe]
Dla 120 produktów, z których 9 jest wadliwych, zapisz dwa modele: z niezależnym losowaniem i z losowaniem bez zwracania ze skończonej partii.
\end{zadanie}
\begin{zadanie}[dowodowe]
Pokaż, że gdy $X\sim N(\theta^2,1)$, parametr $\theta\in\mathbb R$ nie jest identyfikowalny; sformułuj konsekwencję dla zgodności estymacji.
\end{zadanie}
\begin{wskazowka}
Porównaj prawa dla $\theta$ i $-\theta$.
\end{wskazowka}
\begin{zadanie}[praktyczne]
W badaniu pacjentów wykonano cztery pomiary tej samej osoby. Ustal jednostkę niezależną i zaproponuj poprawny model danych.
\end{zadanie}
\begin{wskazowka}
Zastanów się nad korelacją wewnątrz osoby.
\end{wskazowka}
\section{Statystyki i rozkłady z próby}
\label{zadania:IV:2}
\noindent\textit{Teoria:} tom IV, rozdział~\ref{IV-ch:statystyki} (\hyperref[IV-ch:statystyki]{przejdź do materiału}).\par\medskip
\begin{zadanie}[podstawowe]
Wyprowadź $\mathbb E S^2=\sigma^2$ dla niezależnej próby z populacji o skończonym drugim momencie.
\end{zadanie}
\begin{zadanie}[dowodowe]
Dla danych $[10,11,9,12,60]$ oblicz średnią, medianę, wariancję i MAD; wybierz odporną statystykę lokalizacji.
\end{zadanie}
\begin{zadanie}[praktyczne]
Znajdź rozkład maksimum $n$ niezależnych zmiennych jednostajnych na $(0,\theta)$ i wskaż, jak estymować $\theta$.
\end{zadanie}
\begin{wskazowka}
Zacznij od dystrybuanty maksimum.
\end{wskazowka}
\section{Funkcja wiarygodności i informacja zawarta w danych}
\label{zadania:IV:3}
\noindent\textit{Teoria:} tom IV, rozdział~\ref{IV-ch:likelihood} (\hyperref[IV-ch:likelihood]{przejdź do materiału}).\par\medskip
\begin{zadanie}[podstawowe]
Wyznacz funkcję wiarygodności i MLE parametru $p$ przy 9 wadach w próbie 120 niezależnych produktów.
\end{zadanie}
\begin{zadanie}[dowodowe]
Dla $X_i\sim\operatorname{Unif}(0,\theta)$ wyznacz MLE, wskazując, dlaczego standardowy warunek score równy zero zawodzi.
\end{zadanie}
\begin{zadanie}[praktyczne]
W modelu Poissona z ekspozycjami $t_i$ wyprowadź MLE wspólnej intensywności i porównaj ze średnią $N_i/t_i$.
\end{zadanie}
\begin{wskazowka}
Zważ obserwacje długościami ekspozycji.
\end{wskazowka}
\section{Statystyczna teoria decyzji: od danych do działania}
\label{zadania:IV:4}
\noindent\textit{Teoria:} tom IV, rozdział~\ref{IV-ch:decision} (\hyperref[IV-ch:decision]{przejdź do materiału}).\par\medskip
\begin{zadanie}[podstawowe]
Dla estymacji $p$ przy stracie kwadratowej wyprowadź ryzyko estymatora $a\bar X+b$, gdzie $X_i\sim\operatorname{Bern}(p)$.
\end{zadanie}
\begin{zadanie}[dowodowe]
W klasyfikacji choroby porównaj optymalny próg decyzji przy kosztach fałszywego alarmu 1 i przeoczenia 20.
\end{zadanie}
\begin{wskazowka}
Zminimalizuj warunkowe oczekiwane straty.
\end{wskazowka}
\begin{zadanie}[praktyczne]
Podaj problem, w którym estymator obciążony ma mniejszy MSE niż nieobciążony, i udowodnij to rachunkowo.
\end{zadanie}
\section{Konstrukcja i porównywanie estymatorów}
\label{zadania:IV:5}
\noindent\textit{Teoria:} tom IV, rozdział~\ref{IV-ch:estimators} (\hyperref[IV-ch:estimators]{przejdź do materiału}).\par\medskip
\begin{zadanie}[podstawowe]
Porównaj metodę momentów i MLE dla parametru skali rozkładu wykładniczego.
\end{zadanie}
\begin{zadanie}[dowodowe]
Dla próby z rozkładu jednostajnego na $(0,\theta)$ porównaj MSE $X_{(n)}$ oraz $\frac{n+1}{n}X_{(n)}$.
\end{zadanie}
\begin{wskazowka}
Wykorzystaj rozkład maksimum.
\end{wskazowka}
\begin{zadanie}[praktyczne]
Zaproponuj odporny estymator lokalizacji dla próbki z 5\% obserwacji odstających i porównaj go ze średnią w symulacji.
\end{zadanie}
\section{Dostateczność, zupełność i optymalna estymacja nieobciążona}
\label{zadania:IV:6}
\noindent\textit{Teoria:} tom IV, rozdział~\ref{IV-ch:sufficiency} (\hyperref[IV-ch:sufficiency]{przejdź do materiału}).\par\medskip
\begin{zadanie}[podstawowe]
Pokaż dostateczność sumy sukcesów dla parametru Bernoulliego poprzez faktoryzację.
\end{zadanie}
\begin{zadanie}[dowodowe]
Skonstruuj estymator UMVU funkcji $p^2$ z dwóch lub więcej prób Bernoulliego.
\end{zadanie}
\begin{wskazowka}
Zacznij od iloczynu dwóch różnych wskaźników i zastosuj Rao--Blackwella.
\end{wskazowka}
\begin{zadanie}[praktyczne]
Wykaż, dlaczego zupełność statystyki jest istotna w twierdzeniu Lehmanna--Scheffégo.
\end{zadanie}
\begin{wskazowka}
Spróbuj porównać dwa estymatory nieobciążone będące funkcjami tej samej statystyki.
\end{wskazowka}
\section{Informacja Fishera, efektywność i granice dokładności estymacji}
\label{zadania:IV:7}
\noindent\textit{Teoria:} tom IV, rozdział~\ref{IV-ch:fisher-information} (\hyperref[IV-ch:fisher-information]{przejdź do materiału}).\par\medskip
\begin{zadanie}[podstawowe]
Wyprowadź informację Fishera w modelu normalnym $N(\mu,\sigma^2)$ przy znanej wariancji.
\end{zadanie}
\begin{zadanie}[dowodowe]
Oblicz granicę Craméra--Rao dla estymatorów intensywności procesu Poissona obserwowanego przez czas $T$.
\end{zadanie}
\begin{zadanie}[praktyczne]
Wyjaśnij na modelu $\operatorname{Unif}(0,\theta)$, które założenie klasycznej nierówności Craméra--Rao zawodzi.
\end{zadanie}
\section{Asymptotyczna teoria estymacji i niepewność w praktyce}
\label{zadania:IV:8}
\noindent\textit{Teoria:} tom IV, rozdział~\ref{IV-ch:asymptotic-estimation} (\hyperref[IV-ch:asymptotic-estimation]{przejdź do materiału}).\par\medskip
\begin{zadanie}[podstawowe]
Zastosuj metodę delty do $g(\bar X)=\log\bar X$ przy dodatniej średniej i skończonej wariancji.
\end{zadanie}
\begin{zadanie}[dowodowe]
Porównaj dla 7 wad na 100 przedziały Walda i Wilsona, wskazując ich zachowanie blisko zera.
\end{zadanie}
\begin{zadanie}[praktyczne]
Wyprowadź drugi rząd metody delty dla $g(x)=x^2$ przy estymacji parametru $\theta=0$.
\end{zadanie}
\begin{wskazowka}
Pierwsza pochodna w zerze zanika.
\end{wskazowka}
\section{Testowanie hipotez: konstrukcja, moc i interpretacja}
\label{zadania:IV:9}
\noindent\textit{Teoria:} tom IV, rozdział~\ref{IV-ch:testing} (\hyperref[IV-ch:testing]{przejdź do materiału}).\par\medskip
\begin{zadanie}[podstawowe]
W teście hipotezy $H_0:p=0{,}05$ przeciw $H_1:p>0{,}05$ wyznacz obszar krytyczny o poziomie co najwyżej 0,05 dla $n=100$.
\end{zadanie}
\begin{zadanie}[dowodowe]
Udowodnij lemat Neymana--Pearsona dla dwóch prostych hipotez i określ rolę randomizacji na granicy.
\end{zadanie}
\begin{zadanie}[praktyczne]
Oblicz moc testu dla $p=0{,}10$ i wyjaśnij, czym różni się od wartości p przy zaobserwowanym wyniku.
\end{zadanie}
\section{Testy parametryczne: średnie, wariancje, częstości i ANOVA}
\label{zadania:IV:10}
\noindent\textit{Teoria:} tom IV, rozdział~\ref{IV-ch:parametric-tests} (\hyperref[IV-ch:parametric-tests]{przejdź do materiału}).\par\medskip
\begin{zadanie}[podstawowe]
Dla dwóch niezależnych prób o różnych wariancjach wyprowadź statystykę Welcha i uzasadnij wybór zamiast testu t z wariancją łączoną.
\end{zadanie}
\begin{zadanie}[dowodowe]
Dla trzech grup zapisz rozkład sum kwadratów w jednoczynnikowej ANOVA i wyprowadź statystykę $F$.
\end{zadanie}
\begin{zadanie}[praktyczne]
Dla sparowanych pomiarów oblicz test t na różnicach i porównaj z błędnym potraktowaniem obu pomiarów jako niezależnych.
\end{zadanie}
\begin{wskazowka}
Kowariancja pomiarów wpływa na wariancję różnicy.
\end{wskazowka}
\section{Przedziały ufności: konstrukcja, pokrycie i interpretacja}
\label{zadania:IV:11}
\noindent\textit{Teoria:} tom IV, rozdział~\ref{IV-ch:confidence-intervals} (\hyperref[IV-ch:confidence-intervals]{przejdź do materiału}).\par\medskip
\begin{zadanie}[podstawowe]
Dla 7 wad na 100 produktów oblicz 95\% przedział Wilsona oraz dokładny Cloppera--Pearsona.
\end{zadanie}
\begin{zadanie}[dowodowe]
Udowodnij dualność dwustronnego testu i przedziału ufności przy poprawnie dobranych poziomach.
\end{zadanie}
\begin{zadanie}[praktyczne]
Wyprowadź jednostronną dolną granicę ufności dla średniej normalnej o nieznanej wariancji.
\end{zadanie}
\section{Wielokrotne testowanie: FWER, FDR i decyzje po selekcji}
\label{zadania:IV:12}
\noindent\textit{Teoria:} tom IV, rozdział~\ref{IV-ch:multiple-testing} (\hyperref[IV-ch:multiple-testing]{przejdź do materiału}).\par\medskip
\begin{zadanie}[podstawowe]
Dla wartości p $(0{,}001,0{,}009,0{,}03,0{,}2)$ zastosuj korekty Bonferroniego, Holma i BH przy $\alpha=0{,}05$.
\end{zadanie}
\begin{zadanie}[dowodowe]
Udowodnij kontrolę FWER procedury Bonferroniego bez założenia niezależności testów.
\end{zadanie}
\begin{zadanie}[praktyczne]
Wskaż problem badawczy, gdzie FDR jest właściwszym kryterium niż FWER, i uzasadnij to kosztami błędów.
\end{zadanie}
\section{Regresja liniowa: estymacja, diagnostyka i predykcja}
\label{zadania:IV:13}
\noindent\textit{Teoria:} tom IV, rozdział~\ref{IV-ch:linear-regression} (\hyperref[IV-ch:linear-regression]{przejdź do materiału}).\par\medskip
\begin{zadanie}[podstawowe]
Dla modelu $Y=X\beta+\varepsilon$ wyprowadź MNK i jego macierz kowariancji przy homoskedastyczności.
\end{zadanie}
\begin{zadanie}[dowodowe]
Na danych z obserwacją wpływową porównaj MNK przed i po usunięciu punktu; wskaż, dlaczego usunięcie wymaga uzasadnienia merytorycznego.
\end{zadanie}
\begin{zadanie}[praktyczne]
Wyprowadź przedział predykcji nowej obserwacji i wyjaśnij, dlaczego jest szerszy niż przedział średniej odpowiedzi.
\end{zadanie}
\section{Uogólnione modele liniowe: regresja logistyczna i Poissona}
\label{zadania:IV:14}
\noindent\textit{Teoria:} tom IV, rozdział~\ref{IV-ch:glm} (\hyperref[IV-ch:glm]{przejdź do materiału}).\par\medskip
\begin{zadanie}[podstawowe]
W modelu logistycznym wyprowadź interpretację $e^{\beta_1}$ jako ilorazu szans i porównaj z ilorazem ryzyka.
\end{zadanie}
\begin{zadanie}[dowodowe]
W modelu Poissona z offsetem $\log t_i$ wyprowadź przewidywaną liczbę zdarzeń i intensywność na jednostkę ekspozycji.
\end{zadanie}
\begin{zadanie}[praktyczne]
Dla danych nadmiernie rozproszonych porównaj Poissona, quasi-Poissona i ujemny dwumianowy, wskazując, co każdy model zmienia.
\end{zadanie}
\section{Wybór, walidacja i regularizacja modeli statystycznych}
\label{zadania:IV:15}
\noindent\textit{Teoria:} tom IV, rozdział~\ref{IV-ch:model-selection} (\hyperref[IV-ch:model-selection]{przejdź do materiału}).\par\medskip
\begin{zadanie}[podstawowe]
Pokaż na przykładzie wielomianów stopni 1--8, dlaczego minimalizacja błędu treningowego nie wybiera modelu o najlepszej predykcji.
\end{zadanie}
\begin{zadanie}[dowodowe]
Wyprowadź rozwiązanie ridge przy macierzy $X^\top X$ odwracalnej albo dodatnim parametrze kary.
\end{zadanie}
\begin{zadanie}[praktyczne]
Zaprojektuj walidację krzyżową bez przecieku informacji dla selekcji zmiennych i standaryzacji.
\end{zadanie}
\begin{wskazowka}
Wszystkie kroki uczone na danych muszą znajdować się wewnątrz pętli walidacji.
\end{wskazowka}
\section{Wnioskowanie z procesów zależnych: autokorelacja i modele szeregów czasowych}
\label{zadania:IV:16}
\noindent\textit{Teoria:} tom IV, rozdział~\ref{IV-ch:time-series} (\hyperref[IV-ch:time-series]{przejdź do materiału}).\par\medskip
\begin{zadanie}[podstawowe]
Dla $X_t=\phi X_{t-1}+\epsilon_t$ z $|\phi|<1$ wyznacz kowariancję i funkcję autokorelacji.
\end{zadanie}
\begin{zadanie}[dowodowe]
Porównaj błąd standardowy średniej dla próby niezależnej i procesu AR(1) o dodatniej autokorelacji.
\end{zadanie}
\begin{zadanie}[praktyczne]
Zaproponuj procedurę podziału na próbę treningową i testową dla prognoz szeregu czasowego; uzasadnij, dlaczego zwykły losowy podział jest ryzykowny.
\end{zadanie}
\section{Podstawy wnioskowania bayesowskiego: model, prior i posterior}
\label{zadania:IV:17}
\noindent\textit{Teoria:} tom IV, rozdział~\ref{IV-ch:bayes-foundations} (\hyperref[IV-ch:bayes-foundations]{przejdź do materiału}).\par\medskip
\begin{zadanie}[podstawowe]
Dla 7 wad na 100 i prioru $\operatorname{Beta}(2,2)$ wyprowadź posterior i jego średnią.
\end{zadanie}
\begin{zadanie}[dowodowe]
Podaj przykład niewłaściwego prioru prowadzącego do niewłaściwego posterioru.
\end{zadanie}
\begin{zadanie}[praktyczne]
Udowodnij, że estymator Bayesa dla straty kwadratowej jest średnią posterioru, gdy istnieje drugi moment.
\end{zadanie}
\section{Rodziny sprzężone: rachunek posterioru i prognoza}
\label{zadania:IV:18}
\noindent\textit{Teoria:} tom IV, rozdział~\ref{IV-ch:conjugate-families} (\hyperref[IV-ch:conjugate-families]{przejdź do materiału}).\par\medskip
\begin{zadanie}[podstawowe]
Wyprowadź posterior beta--dwumianowy i rozkład predykcyjny liczby sukcesów w kolejnych 20 próbach.
\end{zadanie}
\begin{zadanie}[dowodowe]
W modelu gamma--Poisson uwzględnij różne ekspozycje i oblicz posterior intensywności.
\end{zadanie}
\begin{zadanie}[praktyczne]
Porównaj sprzężony model normalny ze znaną wariancją z modelem o nieznanej wariancji; wskaż, kiedy rozkład predykcyjny jest rozkładem t.
\end{zadanie}
\section{Estymacja i predykcja bayesowska: od posterioru do decyzji}
\label{zadania:IV:19}
\noindent\textit{Teoria:} tom IV, rozdział~\ref{IV-ch:bayes-estimation-prediction} (\hyperref[IV-ch:bayes-estimation-prediction]{przejdź do materiału}).\par\medskip
\begin{zadanie}[podstawowe]
Wyprowadź optymalność mediany posterioru przy stracie bezwzględnej.
\end{zadanie}
\begin{zadanie}[dowodowe]
Dla asymetrycznych kosztów przeszacowania i niedoszacowania wyprowadź właściwy kwantyl posterioru.
\end{zadanie}
\begin{zadanie}[praktyczne]
Dla modelu beta--dwumianowego porównaj 95\% przedział wiarygodności parametru z 95\% przedziałem predykcyjnym liczby sukcesów.
\end{zadanie}
\section{Bayesowskie testowanie hipotez, iloraz Bayesa i paradoksy}
\label{zadania:IV:20}
\noindent\textit{Teoria:} tom IV, rozdział~\ref{IV-ch:bayes-model-testing} (\hyperref[IV-ch:bayes-model-testing]{przejdź do materiału}).\par\medskip
\begin{zadanie}[podstawowe]
Dla testu $H_0:\mu=0$ przeciw normalnemu priorowi alternatywy oblicz iloraz Bayesa i porównaj z wartością p.
\end{zadanie}
\begin{zadanie}[dowodowe]
Wyprowadź przykładową sekwencję $z_n=2$ prowadzącą do zjawiska Jeffreysa--Lindleya.
\end{zadanie}
\begin{zadanie}[praktyczne]
Wyjaśnij, dlaczego prior niewłaściwy nie wyznacza automatycznie sensownego ilorazu Bayesa dla dwóch modeli.
\end{zadanie}
\section{Modele hierarchiczne: częściowe łączenie informacji między grupami}
\label{zadania:IV:21}
\noindent\textit{Teoria:} tom IV, rozdział~\ref{IV-ch:hierarchical-bayes} (\hyperref[IV-ch:hierarchical-bayes]{przejdź do materiału}).\par\medskip
\begin{zadanie}[podstawowe]
W modelu hierarchicznym normalny--normalny wyprowadź wagę partial poolingu i podaj jej granice przy małej i dużej liczebności grupy.
\end{zadanie}
\begin{zadanie}[dowodowe]
Porównaj ryzyko estymacji osobnej i shrinkage w przykładzie trzech grup.
\end{zadanie}
\begin{zadanie}[praktyczne]
Zaproponuj walidację modelu hierarchicznego dla nowej grupy i wyjaśnij, dlaczego losowe rozdzielenie obserwacji wewnątrz grup może zafałszować ocenę.
\end{zadanie}
\section{Asymptotyka bayesowska: koncentracja posterioru i twierdzenie Bernsteina--von Misesa}
\label{zadania:IV:22}
\noindent\textit{Teoria:} tom IV, rozdział~\ref{IV-ch:bayesian-asymptotics} (\hyperref[IV-ch:bayesian-asymptotics]{przejdź do materiału}).\par\medskip
\begin{zadanie}[podstawowe]
Dla próby Bernoulliego z prawdziwym parametrem $p\in(0,1)$ wyprowadź asymptotyczne przybliżenie normalne posterioru beta.
\end{zadanie}
\begin{zadanie}[dowodowe]
Pokaż, że przy prawdziwym $p=0$ poprawna skala posterioru jest rzędu $n^{-1}$, a nie $n^{-1/2}$.
\end{zadanie}
\begin{zadanie}[praktyczne]
Sformułuj problem, w którym multimodalność lub nieidentyfikowalność uniemożliwia bezpośrednie zastosowanie twierdzenia Bernsteina--von Misesa.
\end{zadanie}
\section{Metody Monte Carlo: symulacja, błąd i redukcja wariancji}
\label{zadania:IV:23}
\noindent\textit{Teoria:} tom IV, rozdział~\ref{IV-ch:monte-carlo} (\hyperref[IV-ch:monte-carlo]{przejdź do materiału}).\par\medskip
\begin{zadanie}[podstawowe]
Wyprowadź nieobciążoność estymatora Monte Carlo i jego błąd standardowy.
\end{zadanie}
\begin{zadanie}[dowodowe]
Dla estymacji $\mathbb P(Z>5)$, $Z\sim N(0,1)$, zaproponuj rozkład do importance sampling i poprawne wagi.
\end{zadanie}
\begin{zadanie}[praktyczne]
Wyprowadź optymalny współczynnik zmiennej kontrolnej przy znanej średniej kontrolnej; porównaj wariancje.
\end{zadanie}
\section{Metody MCMC: łańcuchy Markowa, estymacja i diagnostyka}
\label{zadania:IV:24}
\noindent\textit{Teoria:} tom IV, rozdział~\ref{IV-ch:mcmc} (\hyperref[IV-ch:mcmc]{przejdź do materiału}).\par\medskip
\begin{zadanie}[podstawowe]
Dla rozkładu docelowego $\pi(x)\propto e^{-x^2/2}$ skonstruuj algorytm random-walk Metropolis i udowodnij szczegółową równowagę.
\end{zadanie}
\begin{zadanie}[dowodowe]
Porównaj autokorelację i efektywną liczebność próby dla dwóch łańcuchów z różnym krokiem propozycji.
\end{zadanie}
\begin{zadanie}[praktyczne]
Zbuduj przykład dwumodalnego posterioru, w którym cztery łańcuchy startujące w tym samym modzie mają myląco dobre $\widehat R$.
\end{zadanie}
\section{Bootstrap, jackknife i testy permutacyjne: wnioskowanie przez ponowne próbkowanie}
\label{zadania:IV:25}
\noindent\textit{Teoria:} tom IV, rozdział~\ref{IV-ch:resampling} (\hyperref[IV-ch:resampling]{przejdź do materiału}).\par\medskip
\begin{zadanie}[podstawowe]
Wyprowadź dokładną wariancję warunkową bootstrapowej średniej empirycznej.
\end{zadanie}
\begin{zadanie}[dowodowe]
W przypadku 0 sukcesów na 20 prób porównaj naiwny bootstrap przedziału odsetka z przedziałem Wilsona.
\end{zadanie}
\begin{zadanie}[praktyczne]
Dla procesu AR(1) zaproponuj bootstrap blokowy i uzasadnij, dlaczego resampling pojedynczych obserwacji nie zachowuje autokorelacji.
\end{zadanie}
\section{Metody nieparametryczne i odporne: rozkład, rangi i obserwacje odstające}
\label{zadania:IV:26}
\noindent\textit{Teoria:} tom IV, rozdział~\ref{IV-ch:nonparametric} (\hyperref[IV-ch:nonparametric]{przejdź do materiału}).\par\medskip
\begin{zadanie}[podstawowe]
Dla dwóch prób oblicz statystykę Manna--Whitneya i zapisz hipotezę o porównaniu rozkładów.
\end{zadanie}
\begin{zadanie}[dowodowe]
Udowodnij, że mediana minimalizuje sumę modułów odchyleń.
\end{zadanie}
\begin{zadanie}[praktyczne]
Dla próby z pojedynczą dużą obserwacją odstającą porównaj średnią, medianę i estymator Hubera; wyjaśnij wybór funkcji straty.
\end{zadanie}
\section{Diagnostyka rozkładu i założeń: normalność, wariancje i zależność}
\label{zadania:IV:27}
\noindent\textit{Teoria:} tom IV, rozdział~\ref{IV-ch:diagnostics} (\hyperref[IV-ch:diagnostics]{przejdź do materiału}).\par\medskip
\begin{zadanie}[podstawowe]
Wyprowadź niezmienniczość statystyki Shapiro--Wilka względem dodatniego przekształcenia afinicznego danych.
\end{zadanie}
\begin{zadanie}[dowodowe]
Zasymuluj próbki normalne i t-Studenta o małej liczbie stopni swobody; porównaj moc Shapiro--Wilka i Jarque'a--Bery.
\end{zadanie}
\begin{zadanie}[praktyczne]
Dla reszt modelu regresji z heteroskedastycznością porównaj testy Levene'a, Browna--Forsythe'a i Bartletta, wskazując ich ograniczenia.
\end{zadanie}
\section{ANOVA i ANCOVA: modele, interakcje, porównania i zastosowania}
\label{zadania:IV:28}
\noindent\textit{Teoria:} tom IV, rozdział~\ref{IV-ch:anova-ancova} (\hyperref[IV-ch:anova-ancova]{przejdź do materiału}).\par\medskip
\begin{zadanie}[podstawowe]
Dla dwuczynnikowego planu $2\times3$ zapisz model ANOVA z interakcją i wyjaśnij różnice między efektem głównym i prostym.
\end{zadanie}
\begin{zadanie}[dowodowe]
Dla dwóch grup o różnych rozkładach wieku porównaj różnicę średnich surowych i wynik ANCOVA przy wspólnym nachyleniu.
\end{zadanie}
\begin{zadanie}[praktyczne]
Wyjaśnij na niezrównoważonym planie, dlaczego sumy kwadratów typów I, II i III mogą prowadzić do różnych testów.
\end{zadanie}
\section{MANOVA i MANCOVA: wnioskowanie o wielu wynikach jednocześnie}
\label{zadania:IV:29}
\noindent\textit{Teoria:} tom IV, rozdział~\ref{IV-ch:manova-mancova} (\hyperref[IV-ch:manova-mancova]{przejdź do materiału}).\par\medskip
\begin{zadanie}[podstawowe]
Z danych z dwiema zmiennymi wynikowymi skonstruuj macierze $H,E$ i oblicz cztery statystyki MANOVA.
\end{zadanie}
\begin{zadanie}[dowodowe]
Wyprowadź niezmienniczość lambdy Wilksa względem odwracalnej liniowej zmiany współrzędnych zmiennych wynikowych.
\end{zadanie}
\begin{zadanie}[praktyczne]
Porównaj MANOVA z dwoma osobnymi testami t z korektą Holma, wskazując różne testowane hipotezy.
\end{zadanie}
\section{Pomiary powtarzane, sferyczność i modele mieszane}
\label{zadania:IV:30}
\noindent\textit{Teoria:} tom IV, rozdział~\ref{IV-ch:repeated-measures} (\hyperref[IV-ch:repeated-measures]{przejdź do materiału}).\par\medskip
\begin{zadanie}[podstawowe]
Dla czterech pomiarów na osobę zbuduj kontrprzykład niespełnienia sferyczności i oblicz $\epsilon_{GG}$.
\end{zadanie}
\begin{zadanie}[dowodowe]
Porównaj stopnie swobody i wartości p ANOVA powtarzanych pomiarów przed i po korekcie Greenhouse'a--Geissera.
\end{zadanie}
\begin{zadanie}[praktyczne]
Zaproponuj liniowy model mieszany dla niekompletnych wizyt o nieregularnych terminach; opisz założenia dotyczące brakujących danych.
\end{zadanie}
\section{Eksperymenty, A/B testing i wnioskowanie przyczynowe}
\label{zadania:IV:31}
\noindent\textit{Teoria:} tom IV, rozdział~\ref{IV-ch:causal-ab} (\hyperref[IV-ch:causal-ab]{przejdź do materiału}).\par\medskip
\begin{zadanie}[podstawowe]
Dla eksperymentu A/B z losowym przydziałem wyprowadź nieobciążoność różnicy średnich dla średniego efektu przyczynowego.
\end{zadanie}
\begin{zadanie}[dowodowe]
Podaj przykład paradoksu Simpsona przy nierównym udziale grup w warstwach i wyjaśnij rolę standaryzacji.
\end{zadanie}
\begin{zadanie}[praktyczne]
Zapisz warunki identyfikacji ATE za pomocą propensity score i wskaż, dlaczego dodatniość jest niezbędna.
\end{zadanie}
\section{Laboratorium statystyczne: wybór metody, analiza i raportowanie}
\label{zadania:IV:32}
\noindent\textit{Teoria:} tom IV, rozdział~\ref{IV-ch:stat-laboratory} (\hyperref[IV-ch:stat-laboratory]{przejdź do materiału}).\par\medskip
\begin{zadanie}[podstawowe]
Przygotuj analizę 1000 niezależnych wyników A/B: wybierz estymand, test, przedział ufności i miarę efektu; uzasadnij każdy wybór.
\end{zadanie}
\begin{zadanie}[dowodowe]
Zestaw na tych samych danych podejście klasyczne i bayesowskie, wskazując różnicę interpretacyjną.
\end{zadanie}
\begin{zadanie}[praktyczne]
Stwórz protokół raportowania analizy z brakami danych, wieloma testami i obserwacjami odstającymi.
\end{zadanie}
\begin{wskazowka}
Podziel decyzje na podjęte przed zobaczeniem danych i eksploracyjne.
\end{wskazowka}

\end{document}
